\FormSection{1. Thông tin chung:}
\FormLine{Mã số đề tài: \ProjectCode}
\FormLine{Tên đề tài: \ProjectTitle}
\FormLine{Sinh viên chủ nhiệm đề tài: \PrincipalName}
\FormLine{Lớp: DI2396F1 -- Khóa 49; Trường/Khoa/Viện: Trường CNTT\&TT.}
\FormIndented{Năm thứ: \PrincipalYear\qquad Số năm đào tạo: \TrainingYears}
\FormLine{Người hướng dẫn: \SupervisorName}

\FormSection{2. Mục tiêu đề tài:}
Xây dựng hệ thống ứng dụng trí tuệ nhân tạo trên thiết bị di động để tiếp nhận, phân tích thông tin và hỗ trợ ra quyết định cứu hộ bão lũ trong điều kiện mạng không ổn định, hướng tới giảm thiểu thiệt hại về người và tài sản.

\FormSection{3. Tính mới và sáng tạo:}
Đề tài đề xuất cách tiếp cận kết hợp phân tích dữ liệu đa phương thức, trí tuệ nhân tạo tại biên và phân cụm sự kiện cứu hộ trong một hệ thống thống nhất.

Thay vì phụ thuộc hoàn toàn vào máy chủ, hệ thống hướng tới thực hiện suy luận AI trực tiếp trên thiết bị di động, sử dụng mô hình học sâu được tối ưu hóa nhằm giảm yêu cầu tài nguyên và hạn chế truyền tải dữ liệu dung lượng lớn.

Bên cạnh đó, việc kết hợp thông tin vị trí, thời gian và nội dung cứu hộ để phát hiện, gom nhóm các yêu cầu liên quan giúp hạn chế thông tin trùng lặp và hỗ trợ trực quan hóa tình huống cứu hộ trên bản đồ số.

\FormSection{4. Kết quả nghiên cứu:}
Đề tài tập trung nghiên cứu quy trình xây dựng dữ liệu cứu hộ bão lũ, phân tích hình ảnh và văn bản bằng học sâu, tối ưu hóa mô hình AI cho thiết bị di động và thiết kế hệ thống tiếp nhận, phân cụm thông tin cứu hộ.

Đối với bài toán phân loại hình ảnh, mô hình được đánh giá trên tập kiểm thử gồm 256 ảnh. Với cấu hình seed 1024, mô hình đạt Accuracy \NewImageAccuracy\% và macro-F1 \NewImageMacroFOne\%, ghi nhận 3 lỗi nghiêm trọng. Thử nghiệm cắt tỉa mô hình đạt Accuracy \PruningAccuracy\% với \PruningCriticalCount{} lỗi nghiêm trọng.

Các nội dung về phân tích văn bản, phân cụm sự kiện, đồng bộ dữ liệu và hoạt động trong điều kiện mạng yếu được xác định theo mục tiêu nghiên cứu của đề tài. Hiệu quả của từng thành phần cần được đánh giá bằng các kết quả thực nghiệm tương ứng.

Trong thử nghiệm mạng yếu giả lập với băng thông 50 kbit/s và độ trễ khứ hồi 600 ms, chế độ chỉ gửi thông tin cơ bản và chế độ gửi ảnh nén đều thành công ở 257/257 lượt (100\%), trong khi gửi ảnh gốc đạt 230/257 lượt (89,49\%). Kết quả cho thấy hai chế độ gửi ít dữ liệu hơn có tỷ lệ truyền thành công cao hơn gửi ảnh gốc trong điều kiện thử nghiệm. Các phép đo được thực hiện riêng cho từng chế độ trên máy tính; khả năng truyền qua mạng di động thực tế cần được kiểm chứng thêm.

Các kết quả hiện có phản ánh khả năng triển khai và đánh giá nguyên mẫu trong môi trường nghiên cứu, chưa đủ cơ sở để khẳng định hiệu quả điều phối cứu hộ trong điều kiện thiên tai thực tế.

\FormSection{5. Sản phẩm:}
Đề tài đăng ký ba sản phẩm khoa học và công nghệ chính:
\FormIndented{-- Ứng dụng di động hỗ trợ cứu hộ thiên tai, tích hợp mô hình AI và chức năng gửi thông tin cầu cứu.}
\FormIndented{-- Website quản lý thông tin cứu hộ, trực quan hóa vị trí cầu cứu và phân cụm sự kiện trên bản đồ số.}
\FormIndented{-- Bộ mô hình AI đa phương thức được huấn luyện, tinh chỉnh và tối ưu hóa để phục vụ phân loại thông tin cứu hộ.}
\par\noindent Ngoài ra, đề tài đăng ký các sản phẩm theo quy định của Đại học Cần Thơ gồm bản tin, báo cáo tóm tắt và video giới thiệu kết quả nghiên cứu có thời lượng tối đa 02 phút.

\FormSection{6. Công bố khoa học từ kết quả nghiên cứu của đề tài, hoặc nhận xét, đánh giá của cơ sở đã áp dụng các kết quả nghiên cứu (nếu có):}
% Title and author order follow paper/main.tex. Publication status was confirmed
% by the project team on 24/10/2026. CCIS/Scopus Q3 follows the ISDS 2026 CFP:
% https://isds.ctu.edu.vn/2026/author (accessed 24/10/2026).
% The conference year is not a confirmed publication year; volume/year await
% the publisher's announcement. This body is shared with the Vietnamese summary.
\begingroup
\fontsize{11bp}{13bp}\selectfont
\setlength{\tabcolsep}{3bp}
\renewcommand{\arraystretch}{1.15}
\noindent\begin{tabularx}{\linewidth}{|>{\centering\arraybackslash}p{5mm}|>{\raggedright\arraybackslash}X|P{36mm}|P{30mm}|P{17mm}|P{17mm}|}
\hline
\multicolumn{1}{|c|}{\textbf{Stt}} &
\multicolumn{1}{c|}{\textbf{Tên bài báo}} &
\multicolumn{1}{c|}{\shortstack{\strut\textbf{Tác giả/}\\\strut\textbf{nhóm tác giả}}} &
\multicolumn{1}{c|}{\textbf{Tên tạp chí}} &
\multicolumn{1}{c|}{\textbf{Số tạp chí}} &
\multicolumn{1}{c|}{\shortstack{\strut\textbf{Năm}\\\strut\textbf{xuất bản}}} \\
\hline
1. &
Stress-Testing Product-Gated Clustering and Bounded Priority Ranking for Flood-Rescue Reports: A Synthetic Study &
Thanh-Trong Nguyen;\newline Ngoc-Anh Le;\newline Nhu-Quynh Nguyen;\newline Tuong-Hung Cao;\newline Hung-Thinh Ngo;\newline Xuan-Phuong Chau;\newline Lan Phuong Phan;\newline Thanh-Khoa Nguyen &
Communications in Computer and Information Science (CCIS), Springer; kỷ yếu ISDS 2026 (Scopus Q3) &
Chưa công bố &
Chưa công bố \\
\hline
\end{tabularx}\par
\endgroup
\noindent\textit{Bài báo đã được chấp nhận tại ISDS 2026 và đang chuẩn bị xuất bản. CCIS là bộ sách kỷ yếu hội nghị; số tập và năm xuất bản sẽ được cập nhật theo thông tin chính thức của nhà xuất bản.}

\FormSection{7. Đóng góp về mặt kinh tế -- xã hội, giáo dục và đào tạo, an ninh, quốc phòng và khả năng áp dụng của đề tài:}
Về giáo dục và đào tạo, đề tài tạo điều kiện để sinh viên vận dụng kiến thức về học sâu, xử lý dữ liệu đa phương thức, Edge AI, phát triển ứng dụng di động và hệ thống thông tin vào bài toán thực tiễn.

Về khoa học và công nghệ, nghiên cứu góp phần khảo sát khả năng xử lý thông tin cứu hộ tại biên, tối ưu hóa mô hình học sâu và kết hợp phân cụm sự kiện với trực quan hóa dữ liệu không gian.

Về kinh tế -- xã hội, hệ thống có tiềm năng hỗ trợ tiếp nhận, tổng hợp và phân loại thông tin cứu hộ trong điều kiện hạ tầng mạng bị hạn chế, qua đó góp phần nâng cao khả năng hỗ trợ ra quyết định và điều phối nguồn lực khi xảy ra thiên tai.

Các lợi ích thực tế về thời gian cứu hộ, chi phí vận hành và mức độ giảm thiệt hại cần được kiểm chứng trong những nghiên cứu và thử nghiệm triển khai tiếp theo.

\FormSection{8. Hiệu quả, phương thức chuyển giao kết quả nghiên cứu và khả năng áp dụng:}
Các sản phẩm của đề tài dự kiến được đóng gói và chuyển giao cho Trường Công nghệ Thông tin và Truyền thông, Đại học Cần Thơ để phục vụ lưu trữ, giảng dạy, học tập và nghiên cứu.

Kết quả nghiên cứu có thể làm cơ sở để phát triển các ứng dụng hỗ trợ tiếp nhận và quản lý thông tin cứu hộ thiên tai, đặc biệt trong môi trường kết nối mạng hạn chế.

Để triển khai thực tế, hệ thống cần tiếp tục hoàn thiện khả năng phân tích đa phương thức, tối ưu hóa mô hình trên thiết bị di động, kiểm thử trong điều kiện mạng yếu và đánh giá độ tin cậy của quá trình phân cụm, đồng bộ dữ liệu.

Việc ứng dụng vào công tác cứu hộ thực tế cần có thêm thử nghiệm, xác nhận chuyên môn và sự phối hợp với các cơ quan chức năng.
